\section{The four conditions, run on a commission I would rather not count}
\label{sec:F.2}
\ref{sec:F.1} says a commission counts when no party to the conflict controls the body that mandated it, and when its evidence and methodology are published, it offered the parties a reply and recorded the answer or the refusal, a later merits judgment displaces it, and its mandate was fixed before the finding it reaches. The test is run here on the Organisation for the Prohibition of Chemical Weapons' Investigation and Identification Team, which attributes chemical attacks in Syria. Its mandate came from a political majority: the Conference of States Parties adopted it on 27 June 2018 by 82 votes to 24, over the protest of Russia, which said the move went beyond the organization's mandate \autocite{opcw2018attribution,bbc2018opcwblame}. The upstream investigation of one incident it later attributed, Douma in April 2018, was the subject of leaked internal dissent \autocite{democracynow2019douma,opcw2020breaches}. A body with that origin and that history is one I would rather not count.

\emph{The threshold.} Met. The mandate was adopted by vote over the protest of Russia, Syria's principal ally, so no party to the conflict controlled the decision; the vote record is what shows it.

\emph{Evidence and methodology published.} Met. Its five reports are public in the organization's six languages, and its methods, its rules on samples and chain of custody, and its standard of proof, “reasonable grounds”, the standard fact-finding missions and commissions of inquiry use, are published \autocite{opcw2026iitpage}. The Douma dissent concerned the upstream fact-finding mission, a different body whose mandate excludes attribution \autocite{opcwffm}; the organization's own inquiry found the leaked assessment to be a personal document prepared on incomplete information \autocite{opcw2020breaches}, and an outside analysis reached the same conclusion on the engineering \autocite{bellingcat2020henderson}. Both are accounts by parties to the dispute, and they go on the record beside the finding, as the rule above requires.

\emph{A reply offered and recorded.} Met in substance. Syria's former government didn't cooperate with any investigation before 2025, and the organization records that it considered that government's positions and found no concrete information supporting them \autocite{opcw2026iitkafrzeita}; Syria's and Russia's statements on the first report are published beside everyone else's. Whether draft findings were sent for reply before each report is in the reports, which this translation didn't read.

\emph{Displaced by a later merits judgment.} Met in design and never exercised. The team disclaims judicial status and individual responsibility \autocite{opcw2026iitpage}, and no court has reached the merits of any incident it attributed: the French, German and Swedish proceedings concern other attacks, and the Court's case against Syria is under the Torture Convention \autocite{osji2021frenchcw,icj2023canadasyria}.

\emph{Mandate fixed before the finding.} Met on the wording: the mandate dates from June 2018 and the first finding from April 2020 \autocite{opcw2020iitltamenah}. It is contestable in spirit, since the mandate reaches only Syria and was adopted after most of the incidents it would examine. But a rule requiring the mandate to precede the events would fail nearly every commission of inquiry, since they are created in response to events, and the mandate's generic terms have produced a finding against a non-state group as well as against the state \autocite{opcw2024iitmarea}.

On the four conditions, the team counts. Its Douma attribution \autocite{opcw2023iitdouma} is a finding of the second kind, displaced by any merits judgment that comes, and its political origin doesn't change that. Counting it licenses nothing under the force term, which lifts only for the exceptions \S\ref{sec:6.1} names; what it settles is that my dislike of a body is not one of the four conditions.
